\subsection{Normal Sampling Pivots}

For normal samples, the sample mean and sample variance separate exactly.  This is the finite-sample source of the classical $z$, $\chi^2$, and Student $t$ pivots used in normal-theory confidence intervals and tests.

\begin{definition} \label{def:normal-sample-mean-variance} \hlt{Normal Sample Statistics}
\begin{enumerate}[label=(\roman*)]
\setlength{\itemsep}{0pt}
\item \hlt{Sample}: $X_1,\ldots,X_n$ are independent with $X_i\sim N(\mu,\sigma^2)$, $n\geq2$, and $\sigma^2>0$.
\item \hlt{Sample mean}: $\overline{X}\defeq n^{-1}\sum\limits_{i=1}^nX_i$ estimates the location $\mu$.
\item \hlt{Sample variance}: $S^2\defeq(n-1)^{-1}\sum\limits_{i=1}^n(X_i-\overline{X})^2$ estimates the scale $\sigma^2$.
\end{enumerate}
\end{definition}

\begin{proposition} \label{prop:normal-sample-mean-law} \hlt{Normal Sample Mean}\\
Let $X_1,\ldots,X_n$ be independent with $X_i\sim N(\mu,\sigma^2)$, and set $\overline{X}\defeq n^{-1}\sum\limits_{i=1}^nX_i$. Then $\overline{X}\sim N(\mu,\frac{\sigma^2}{n})$.
\begin{proof} Since independent normal sums are normal, $\overline{X}$ is normal.\\
Also $\mathbb{E}\overline{X}=n^{-1}\sum\limits_i\mathbb{E}X_i=\mu$ and $\operatorname{Var}(\overline{X})=n^{-2}\sum\limits_i\operatorname{Var}(X_i)=\frac{\sigma^2}{n}$. Thus $\overline{X}\sim N(\mu,\frac{\sigma^2}{n})$.\\
\hlt{Structural Commentary}: Averaging preserves normality and divides variance by $n$.\\
\hlt{Key Insight}: The sample mean is an exact normal pivot center, not an asymptotic approximation.
\end{proof}
\end{proposition}

\begin{theorem} \label{thm:normal-mean-residual-independence} \hlt{Mean-Residual Independence}\\
Let $X_1,\ldots,X_n$ be independent with $X_i\sim N(\mu,\sigma^2)$, set $\overline{X}\defeq n^{-1}\sum\limits_iX_i$, $R\defeq(X_1-\overline{X},\ldots,X_n-\overline{X})$,\\
and $S^2\defeq(n-1)^{-1}\sum\limits_iR_i^2$. Then $\overline{X}\perp\!\!\!\perp R$, hence $\overline{X}\perp\!\!\!\perp S^2$.
\begin{proof} Let $Z_i\defeq\frac{X_i-\mu}{\sigma}$, so $Z\sim N_n(0,I_n)$. Choose orthogonal $Q$ with first row $n^{-\frac{1}{2}}\mathbf{1}^\top$, set $Y\defeq QZ$.\\
Then $Y\sim N_n(0,I_n)$, $Y_1\perp\!\!\!\perp(Y_2,\ldots,Y_n)$. Since $\overline{X}=\mu+\frac{\sigma Y_1}{\sqrt n}$ and $R=\sigma(I-n^{-1}\mathbf{1}\mathbf{1}^\top)Z=\sigma Q^\top(0,Y_2,\ldots,Y_n)^\top$, $\overline{X}$ is function of $Y_1$, $R$ is function of $(Y_2,\ldots,Y_n)$. Therefore $\overline{X}\perp\!\!\!\perp R$, by $S^2=(n-1)^{-1}\sum\limits_iR_i^2$, $\overline{X}\perp\!\!\!\perp S^2$.\\
\hlt{Structural Commentary}: Mean and residuals are orthogonal projections.\\
\hlt{Key Insight}: For Gaussian vectors, orthogonal components are independent.
\end{proof}
\end{theorem}

\begin{theorem} \label{thm:normal-sample-variance-chi-square} \hlt{Sample Variance Chi-Square Law}\\
Given independent $X_1,\ldots,X_n$; $X_i\sim N(\mu,\sigma^2)$;$S^2\defeq(n-1)^{-1}\sum\limits_{i=1}^n(X_i-\overline{X})^2$. Then $\frac{(n-1)S^2}{\sigma^2}\sim\chi^2_{n-1}$.
\begin{proof} Let $Z_i\defeq\frac{X_i-\mu}{\sigma}$, so $Z=(Z_1,\ldots,Z_n)\sim N_n(0,I_n)$.\\
Choose orthogonal matrix $Q$ with first row $n^{-\frac{1}{2}}(1,\ldots,1)$, set $Y\defeq QZ$.\\
Orthogonal invariance gives $Y\sim N_n(0,I_n)$ with independent coordinates.\\
Then $Y_1=\frac{\sqrt n(\overline{X}-\mu)}{\sigma}$ and $\sum\limits_{j=2}^nY_j^2=\sum\limits_{i=1}^n(Z_i-\overline{Z})^2=\frac{(n-1)S^2}{\sigma^2}$.\\
Thus $\frac{(n-1)S^2}{\sigma^2}$ is the sum of squares of $n-1$ independent $N(0,1)$ variables, hence $\chi^2_{n-1}$.\\
\hlt{Structural Commentary}: Cantering leaves $n-1$ residual directions.\\
\hlt{Key Insight}: The sample variance is the squared Gaussian residual length scaled by $\sigma^2$.
\end{proof}
\end{theorem}

\begin{corollary} \label{cor:normal-student-t-pivot} \hlt{Student Pivot for a Normal Mean}\\
Let $X_1,\ldots,X_n$ be independent; $X_i\sim N(\mu,\sigma^2)$, $\overline{X}\defeq n^{-1}\sum\limits_iX_i$, $S^2\defeq(n-1)^{-1}\sum\limits_i(X_i-\overline{X})^2$; $\frac{\overline{X}-\mu}{\frac{S}{\sqrt n}}\sim t_{n-1}$.
\begin{proof} Let $Z\defeq\frac{\sqrt n(\overline{X}-\mu)}{\sigma}$ and $V\defeq\frac{(n-1)S^2}{\sigma^2}$. By normal mean law, $Z\sim N(0,1)$. By sample-variance chi-square law, $V\sim\chi^2_{n-1}$. By mean-residual independence, $Z\perp\!\!\!\perp V$. Therefore $\frac{\overline{X}-\mu}{\frac{S}{\sqrt n}}=\frac{Z}{\sqrt{\frac{V}{n-1}}}\sim t_{n-1}$.\\
\hlt{Structural Commentary}: Replacing $\sigma$ by $S$ yields a Student pivot.\\
\hlt{Key Insight}: Student's $t$ is a normal variable divided by an independent chi-square scale estimate.
\end{proof}
\end{corollary}

\begin{method} \label{meth:normal-mean-exact-confidence-intervals} \hlt{Exact Normal Mean Intervals}\\
Let $X_1,\ldots,X_n$ be independent $N(\mu,\sigma^2)$ variables, with $\overline{X}\defeq n^{-1}\sum\limits_iX_i$ and $S^2\defeq(n-1)^{-1}\sum\limits_i(X_i-\overline{X})^2$.\\
Below cases have exact coverage $1-\alpha$.
\begin{enumerate}[label=(\roman*)]
\setlength{\itemsep}{0pt}
\item \hlt{Known variance}: if $\sigma$ is known, $C_\mu(X)=[\overline{X}-\frac{z_{1-\frac{\alpha}{2}}\sigma}{\sqrt n},\overline{X}+\frac{z_{1-\frac{\alpha}{2}}\sigma}{\sqrt n}]$
\item \hlt{Unknown variance}: if $\sigma$ is unknown, $C_\mu(X)=[\overline{X}-\frac{t_{n-1,1-\frac{\alpha}{2}}S}{\sqrt n},\overline{X}+\frac{t_{n-1,1-\frac{\alpha}{2}}S}{\sqrt n}]$
\item \hlt{Variance}: $C_{\sigma^2}(X)=[\frac{(n-1)S^2}{\chi^2_{n-1,1-\frac{\alpha}{2}}},\frac{(n-1)S^2}{\chi^2_{n-1,\frac{\alpha}{2}}}]$
\end{enumerate}
\end{method}
